\statement{C.44}{Proved}{A full unit and relative pronilpotent conjugacy for the regular pants action}
Let $F=\widehat F(x,y)$, put $z=(xy)^{-1}$, and use the free
basis $(x,z)$. Let $\sigma$ be a regular negative pants involution
as in C.43, normalized so that $\sigma(z)=z^{-1}$ literally.
Put $K=\langle\!\langle z\rangle\!\rangle_F$ and
$A=F/K=\widehat{\mathbb Z}$, with generator the image of $x$.
There is a unique $g\in F$ such that
\[
                    \sigma(g)g^{-1}=z.                       \tag{C.44.1}
\]
Its image $a\in A$ belongs to $\widehat{\mathbb Z}^{\times}$.
Moreover, the substitution $\varphi(x)=g$, $\varphi(z)=z$
induces an automorphism of
\[
 G=F/R_{\mathrm{nil}}(K),\qquad
 R_{\mathrm{nil}}(K)=\bigcap_p O^p(K),
                                                               \tag{C.44.2}
\]
where $O^p(K)$ is the residual for finite $p$-group quotients,
and conjugates the standard regular action
$\rho(x)=zx$, $\rho(z)=z^{-1}$ to $\sigma$ on $G$.
The conjugacy asserted here is in $\operatorname{Aut}(G)$.
No preservation of the original three peripheral classes by
$\varphi$, or generation of the full group $F$ by $g,z$, is asserted.

\begin{proof}
The regular action is the identity on $A$: on the original
abelianization it is negative interchange, and killing $z$ imposes
$y=x^{-1}$. In $F\rtimes_\sigma C_2$, the sections $\sigma$ and
$z\sigma$ are involutions. C.43 makes them $F$-conjugate, and
applying $\sigma$ to a conjugator gives (C.44.1) with the stated
convention. The ratio of two such solutions is fixed by $\sigma$;
C.43 therefore gives uniqueness. The literal normalization is
available by choosing an involutory lift in the local normalizer
of $\overline{\langle z\rangle}$, as in C.43.

Write $\Lambda=\mathbb F_2[[A]]$ and let $T$ denote the image of
$x$. The exact compact Alexander module is
\[
 K/\Phi_2(K)\cong\Lambda[z],\qquad
 \Phi_2(K)=\overline{K^2[K,K]}.                              \tag{C.44.3}
\]
To verify this with the full procyclic head, let
$U_m=\ker(F\to\mathbb Z/m)$. Its Schreier basis is
$x^m$ together with $x^izx^{-i}$, $0\le i<m$.
For $m\mid m'$ the inclusion-induced map on mod-two homology
multiplies the first coordinate by $m'/m$ and sends each boundary
coordinate to its index reduced modulo $m$. In the inverse limit
the first coordinate vanishes, by passage from $m$ to $2m$;
the remaining transitions define precisely $\Lambda$.

This inverse limit is indeed the left side of (C.44.3).
Every continuous character $K\to\mathbb F_2$ extends to an
open neighborhood $KV$: choose $V\trianglelefteq_oF$ with
$K\cap V$ in its kernel and assign value zero to $V$.
This is well-defined because $[V,K]\subseteq V\cap K$.
Every open neighborhood of $K$ contains some $U_m$.
Thus the character group of $K$ is the direct limit of the
character groups of $U_m$, and compact duality gives the asserted
inverse limit. Conjugation by $x$ shifts the boundary coordinates,
proving the claimed module action. Since $\sigma$ is the identity
on $A$, its action on this module is $\Lambda$-linear; since it
inverts $z$, that action is the identity in characteristic two.

Multiplication by $T-1$ is a topological isomorphism from
$\Lambda$ onto its augmentation ideal. Its compact image projects
onto the augmentation ideal in every $\mathbb F_2[\mathbb Z/m]$,
which proves surjectivity. An element of its kernel has constant
coefficients in each such cyclic group ring. The transition from
$2m$ to $m$ sends that constant vector to zero, so all these
coefficients vanish. This proves injectivity; compactness gives
continuity of the inverse. Consequently, for every
$a\in\widehat{\mathbb Z}$ there is a unique continuously varying
$D_a\in\Lambda$ such that
\[
                    (T-1)D_a=T^a-1.                         \tag{C.44.4}
\]
For nonnegative integers this is the usual geometric sum, and its
augmentation is $a$ modulo two; both assertions extend by continuity.

Write $\sigma(x)=\kappa x$, let $\theta=[\kappa]\in\Lambda$,
and write $g=kx^a$ with $k\in K$. The crossed-homomorphism identity,
first on integer powers and then by continuity, gives
$[\sigma(x^a)x^{-a}]=\theta D_a$.
The contribution of $k$ cancels in (C.44.3). Equation (C.44.1)
therefore implies
\[
                           \theta D_a=1.                    \tag{C.44.5}
\]
In particular $D_a$ is a unit, and augmentation makes $a$ odd.
For an odd prime $q$, evaluate $T$ at a primitive $q$th root
$\zeta$ in a finite field of characteristic two. This is a
continuous ring homomorphism factoring through
$\mathbb F_2[\mathbb Z/q]$. If $a=0\pmod q$, (C.44.4) would
give $D_a(\zeta)=0$, contradicting its invertibility. Thus $a$
is nonzero modulo every prime, proving
$a\in\widehat{\mathbb Z}^{\times}$.

For every prime $p$, the same cover calculation gives
$K/\Phi_p(K)=\mathbb F_p[[A]][z]$: pass from $m$ to $pm$ to
kill the $x^m$ coordinate. Put $G_p=F/O^p(K)$,
$P_p=K/O^p(K)$, and let $H_p$ be the closed subgroup generated by
$g,z$ in $G_p$. Its image in $A$ is all of $A$, because $a$ is a
full unit. The elements $g^t z g^{-t}$,
$t\in\widehat{\mathbb Z}$, belong to $H_p\cap P_p$ and have
Frattini images $T^{at}[z]$. Their closed linear span is the
entire module. Hence $H_p\cap P_p$ maps onto $P_p/\Phi(P_p)$
and equals $P_p$ by the pro-$p$ Frattini property. This does not
require finite generation of $P_p$: a proper closed subgroup lies
in a proper open subgroup, and a finite $p$-group quotient then
places it in a maximal subgroup of index $p$. Therefore $H_p=G_p$.

The maximal pronilpotent quotient of $K$ is
$P=K/R_{\mathrm{nil}}(K)=\prod_pP_p$. Indeed the map onto any
finite product of distinct-prime pro-$p$ quotients is surjective,
by the absence of a nontrivial common finite quotient; compactness
then gives the full product. This also identifies the displayed
residual with the intersection of the closed lower central series,
by testing finite nilpotent images. Let $H$ be the closed subgroup
of $G$ generated by $g,z$. For each $p$, its surjection to $G_p$
makes $H\cap P$ surject onto $P_p$: a preimage of an element of
$P_p$ already has zero image in the common quotient $A$.
A closed subgroup of a product of distinct-prime Sylows which
surjects onto each factor is the whole product, by the same finite
product argument and compactness. Thus $H=G$.

The substitution $\varphi$ preserves $K$ and every $O^p(K)$:
precomposing any finite $p$-group quotient of $K$ with
$\varphi|_K$ is again such a quotient. It consequently induces
surjective endomorphisms of $G$ and all $G_p$. A finitely generated
profinite group is Hopfian: pullback under a surjection permutes
its finite set of open normal subgroups of each bounded index,
so the induced maps on their finite intersection quotients are
automorphisms and the kernel is trivial.
Finally (C.44.1) and $\sigma(z)=z^{-1}$ give
$\varphi\rho=\sigma\varphi$ on the two generators. This proves
the stated relative conjugacy.
\end{proof}
